% HAND-WRITTEN fragment -- NOT generated by maestro2tex. Input directly from
% AnalogChapter.tex, after the block sections: it is a system statement built
% on all of them and must follow every section it cites.
%
% Sources of the statements below, so they can be re-checked:
%   A1 offset / RE tracking      tab_ampab_mc_tracking, tab_ampab_mc_offset,
%                                tab_px_ctrl_tracking, tab_px_mc_rezero
%   A2 offset / zero current     tab_tiaab_mc_offset, tab_px_mc_izero,
%                                tab_px_mc_izero_allgain
%   reference DAC                tab_dacr2r12_mc_ideal, tab_dacr2r12_mc_reference,
%                                note_dacr2r12_supply (droop fit: R^2 = 94 %,
%                                residual 931 uV rms against 25.70 mV)
%   converter                    tab_adc_transfer (effective LSB 1.591 mV)
%   feedback resistor            tab_tiarprog_corners, tab_tiarprog_spec,
%                                tab_tiarprog_mc
%   bias generator               tab_biasgen_trim, tab_biasgen_tempco;
%                                MC sigma/mu 12.5 % from Plan.4.Run.2 (bg_mc)
%   channel wiring as built      skill/castalia_anatop_build.il, buildChannel():
%                                afe_out -> SAR Analog_Input directly, no mux,
%                                no switches in the CE/RE/WE paths
%   register file for the switch controls  include/CqAnalog.tex (SWM word)
%   flash / boot ownership       MULTICORE-intro-castalia-2026-07.tex
% Yield figures against the 4-LSB budget are computed here from the published
% mean and sigma of the tables named above, not transcribed from a run.
% Library cell names are deliberately kept OUT of the rendered prose (owner
% request 2026-08-21); they live in these comments instead.

\subsection{Per-Chip Calibration}\label{ss:calibration}

This section documents \emph{design intent}. No calibration has been executed,
in simulation or on silicon; what follows is the procedure the preceding block
characterisations require, the signal paths that procedure needs, and the
residual error it is expected to leave. Paths that do not exist in the current
schematics are marked as design requirements where they arise.

A channel carries two measurement axes and they fail independently. The
\emph{potential axis} polarises the cell to \(V_{\mathrm{CM}} -
v_{\mathrm{ref}}\): the reference DAC of Section~\ref{ss:dacr2r12} sets
\(v_{\mathrm{ref}}\), control amplifier A1 (Section~\ref{ss:ampab}) holds the
reference electrode there, and the shared common-mode DAC sets the working
electrode. The \emph{current axis} reads \(i_{\mathrm{we}} = (v_{\mathrm{out}}
- V_{\mathrm{CM}})/R_f\): transimpedance amplifier A2 converts through the
programmable resistor of Section~\ref{ss:tiarprog} and the converter of
Section~\ref{ss:adc} digitises the result offset-binary about the code at
\(V_{\mathrm{CM}}\).

Neither axis meets a four-code budget uncalibrated. Against
\(\pm4\,\mathrm{LSB}\) of the converter's measured \SI{1.591}{\milli\volt} step
(\SI{6.36}{\milli\volt}), A1's input offset --- mean \SI{0.583}{\milli\volt},
\(\sigma = \SI{5.201}{\milli\volt}\), Table~\ref{tab:ampab-mc-offset} --- passes
\SI{77.6}{\percent} of channels and \SI{36.3}{\percent} of four-channel parts;
carried to the reference electrode as a tracking error
(Table~\ref{tab:ampab-mc-tracking}) it passes \SI{73.2}{\percent} and
\SI{28.6}{\percent}. The current axis is worse by two orders of magnitude,
because the same loop offset is divided by the electrode impedance and
multiplied by \(R_f\): the zero-current offset is \(\sigma = 12.1\), 43.5, 95.7
and \SI{167.3}{LSB} at the 35, 140, 315 and \SI{560}{\kilo\ohm} taps
(Section~\ref{sss:pixel-mc}), for \SIrange{1.9}{25.9}{\percent} of channels
inside \(\pm4\,\mathrm{LSB}\) and effectively no four-channel parts at all.
Both distributions are static offsets, and a static offset is exactly what a
per-chip constant removes.

One property of the chip decides the shape of the whole procedure: \textbf{there
is no on-chip reference of either kind}. The reference DAC's full scale is
whatever its own supply block delivers (\SI{2.166}{\volt} mean, \(\sigma =
\SI{56.8}{\milli\volt}\), Table~\ref{tab:dacr2r12-mc-reference}); the
converter's is the \SI{2.5}{\volt} analog rail; the feedback resistor is poly,
\SIrange{13.3}{26.4}{\kilo\ohm} at its bottom code across corners
(Table~\ref{tab:tiarprog-corners}). None of the three is trimmed, and no two of
them are the same quantity. Absolute accuracy therefore rests on exactly two
external measurements per chip --- one voltage and one current, taken by the
production tester --- and everything else is a ratio the chip measures against
itself. That split is what makes a power-up recalibration possible without
equipment: the terms that drift are ratios and zeros, and those the channel can
re-measure through its own converter.
